\section{Model Training}
\label{sec:training}

The training process for HY-Motion DiT consists of three stages: large-scale pretraining, high-quality fine-tuning, and reinforcement learning. 
To reconcile the trade-off between generalization (requiring massive scale) and precision (requiring high-quality data), we implement a ``scale-then-refine'' curriculum for the supervised training phases (\cref{subsec:stage1} and \cref{subsec:stage2}).
We begin with large-scale pretraining on an expansive dataset comprising over $3,000$ hours of motion data, enabling the model to assimilate a broad motion prior and learn a generalized representation of human movement.
This is followed by a high-quality fine-tuning stage on nearly $400$ hours of meticulously curated and annotated data, which serves to sharpen the model's capabilities and refine the details of the generation.
While these supervised stages effectively model the underlying data distribution, strictly imitating the training distribution does not necessarily maximize perceptual quality or semantic responsiveness.
Therefore, we introduce a final reinforcement learning phase (\cref{subsec:stage3}) to bridge the gap between ``statistical likelihood'' and ``human preference''.
By incorporating feedback from reward models, this stage fine-tunes the model to optimize for perceptual motion quality and maximize adherence to complex user instructions—metrics that are often difficult to capture through supervised loss functions alone.



\subsection{Large-Scale Pretraining}
\label{subsec:stage1}

The primary objective of this foundational stage is general motion prior acquisition.
To maximize the model's exposure to diverse semantic scenarios and kinematic patterns, we train the model on our full motion dataset $\mathcal{D}_{\text{all}}$ consisting of over $3,000$ hours of motion data. 
This dataset prioritizes coverage over precision, comprising both high-quality captures and noisy, in-the-wild video extractions, accompanied by text descriptions ranging from manual annotations to VLM-generated annotations.
In this phase, we employ a standard Flow Matching objective (\cref{eq:flow_matching}) with a constant learning rate $\eta_{\text{pre}}$, aiming to rapidly establish a broad support for the motion distribution.

From a distributional perspective, this stage allows the model to effectively ``learn to move''.
We observe a rapid convergence in training loss, indicating that the model successfully captures the fundamental dynamics of human movement and establishes a robust mapping between diverse textual prompts and plausible poses.
The model demonstrates strong semantic generalization, capable of responding to a vast vocabulary of actions.
However, the generated motion quality inevitably mirrors the mixed nature of the input distribution.
Since a significant portion of $\mathcal{D}_{\text{all}}$ contains noisy data, the model's output distribution remains high-entropy and loosely constrained.
Consequently, while semantically correct, the generated sequences often exhibit artifacts such as high-frequency jitter, foot sliding, and minor anatomical inconsistencies.
We treat this result as an expected outcome: the model has successfully learned a generalized but coarse prior, providing a semantic-rich initialization that is ready to be "sharpened" in the subsequent fine-tuning stage.

\subsection{High-Quality Fine-Tuning}
\label{subsec:stage2}

Following the broad acquisition of motion priors, the second stage focuses on distribution refinement.
To bridge the gap between ``plausible'' and ``high-quality'', we transition the training source to the meticulously curated subset $\mathcal{D}_{\text{HQ}}$.
This dataset is characterized by high-quality, rigorously filtered motion data, paired with manually corrected textual descriptions.
The objective here is to concentrate the model's probability density around the optimal modes of the motion manifold.
Crucially, to refine fine-grained kinematic details while preventing the forgetting of the broad semantic priors learned in \cref{subsec:stage1}, we decay the learning rate to $\eta_{\text{ft}} = 0.1 \times \eta_{\text{pre}}$. 
This conservative optimization strategy ensures that the model polishes its output distribution without collapsing its generative diversity.

Empirically, this phase drives a qualitative evolution from ``roughly correct'' to ``precise and smooth''.
Kinematically, we observe that the model effectively suppresses the noise patterns inherited from the pretraining phase, drastically reducing high-frequency jitter and foot sliding artifacts while enforcing stricter anatomical consistency.
Semantically, the alignment becomes significantly more rigorous; the model demonstrates enhanced sensitivity to fine-grained anatomical instructions—such as accurately distinguishing between ``waving the left hand'' and ``waving the right hand''—which were often ambiguous in the noisy pretraining data.
Importantly, this sharpening of quality and control is achieved without a significant degradation in motion diversity, validating the efficacy of our coarse-to-fine curriculum.


\subsection{Reinforcement Learning}
\label{subsec:stage3}

While supervised learning establishes a strong kinematic prior, optimizing for data likelihood does not strictly equate to maximizing perceptual fidelity or instruction adherence.
Generative models often suffer from a misalignment gap where statistically probable motions may still exhibit physical artifacts or semantic ambiguity.
To bridge this gap, we propose a synergistic two-phase reinforcement learning, progressing from human preference learning to explicit objective constraint satisfaction.

\paragraph{Human Preference Alignment via DPO.}
We first aim to internalize the nuanced standards of human perception.
Since defining a closed-form reward function for complex semantic alignment and motion quality is intractable, we leverage Direct Preference Optimization (DPO)~\citep{Rafailov2023} to steer the policy directly from human feedback.
We utilize the Stage-2 model to generate candidate pairs for a diverse prompt set.
From an annotated pool of $40,000$ pairs, we curate a high-information subset of $9,228$ pairs where human judges identified a distinct ``winner''($\bm{x}_w$) and ``loser''($\bm{x}_l$) based on instruction adherence and visual plausibility.
By maximizing the likelihood margin between these pairs, DPO implicitly models the latent reward function of human perception:
\begin{equation}
    \mathcal{L}_{\text{DPO}}(\pi_{\bm{\theta}};\pi_{\text{ref}}) = - \mathbb{E}_{(\bm{c}, \bm{x}_w, \bm{x}_l) \sim \mathcal{D}_{\text{pref}}} \left[ \log \sigma \left( \beta \log \frac{\pi_{\bm{\theta}}(\bm{x}_w|\bm{c})}{\pi_{\text{ref}}(\bm{x}_w|\bm{c})} - \beta \log \frac{\pi_{\bm{\theta}}(\bm{x}_l|\bm{c})}{\pi_{\text{ref}}(\bm{x}_l|\bm{c})} \right) \right],
\end{equation}
where $\pi_{\text{ref}}$ represents the frozen reference policy, and $\mathcal{D}_{\text{pref}}=\{(\bm{c}^i, \bm{x}_w^i, \bm{x}_l^i)\}_{i=1}^N$.
This phase effectively acts as a semantic anchor, drastically improving the model's ``Pass Rate'' by pruning low-quality modes that, while kinetically valid, fail to align with user intent.

\paragraph{Physics \& Semantic Refinement via Flow-GRPO.}
Building upon the preference-aligned policy, we address the limitations of DPO in enforcing strict physical and semantic boundaries.
We employ Flow-GRPO~\citep{Liu2025}, a variant of Group Relative Policy Optimization tailored for flow matching models, to optimize explicit objectives.
Unlike standard PPO which requires a value network, Flow-GRPO stabilizes training by normalizing advantages within a group of sampled outputs for a given prompt.
We maximize the following objective:
\begin{equation}
    \mathcal{L}_{\text{Flow-GRPO}}(\pi_{\bm{\theta}};\pi_{\text{ref}}) = \mathbb{E}_{\bm{c} \sim \mathcal{D}_\text{GRPO}, \{\bm{x}^i\}_{i=1}^G \sim \pi_{\text{ref}}(\cdot| \bm{c})} \left[ f(r,\hat{A},\bm{\theta}, \epsilon,\beta) \right],
\end{equation}
where the objective function $f$ is averaged over time steps $T$:
\begin{equation}
    \begin{aligned}
        f(r,\hat{A},\bm{\theta}, \epsilon,\beta) &= \frac{1}{G}\sum_{i=1}^G \frac{1}{T} \sum_{t=1}^T \Big( \min \big( r_t^i(\bm{\theta})\hat{A}^i, \text{clip}(r_t^i(\bm{\theta}), 1-\epsilon, 1+\epsilon)\hat{A}^i \big) \\
        & \qquad - \beta D_{\text{KL}}(\pi_{\bm{\theta}} \| \pi_{\text{ref}}) \Big), \
        \text{with}\quad r_t^i(\bm{\theta}) = \frac{p_{\bm{\theta}}(\bm{x}_{t-1}^i|\bm{x}_{t}^i,\bm{c})}{p_{\text{ref}}(\bm{x}_{t-1}^i|\bm{x}_{t}^i,\bm{c})}.
    \end{aligned}
\end{equation}
Here, $\mathcal{D}_\text{GRPO}$ denotes the prompt dataset used for GRPO, $G$ denotes the group size, and $\hat{A}^i$ is the standardized advantage derived from the composite reward $R(\bm{x}^i, \bm{c})$.
Specifically, the reward function integrates two critical metrics:
a semantic reward ($R_{\text{sem}}$)  evaluated by a custom-trained Text-Motion Retrieval (TMR) model \cite{petrovich2023tmr}, and a physical reward ($R_{\text{phy}}$) that imposes hard penalties on artifacts such as foot sliding and root drift.
By explicitly optimizing this objective, the model is fine-tuned to satisfy rigorous kinematic constraints while maximizing semantic precision.
